% \vspace{-1mm}
\section{Conclusion and Discussion} \label{sec:discussion}

In this work, we study privacy leakage when learning models over a mixture of tasks in multitask learning. In particular, we focus on the setting where MTL is used to jointly train a model for many tasks at once by learning a shared representation that captures common features between tasks. We propose a novel, purely black-box \emph{task-inference} threat model, where the adversary's goal is to infer the inclusion of a target task in training given only query access to the smallest shared component when learning jointly over many tasks, the shared representation. By analyzing our task-inference threat model in the context of tracing attacks on Gaussian mean estimation, we find a separation between \textbf{strong} adversaries with access to training samples and \textbf{weak} adversaries with access to fresh samples from the target task's distribution that did not appear in the training dataset. To verify our analysis, we propose purely black-box attacks on machine learning models trained with MTL and we conduct extensive experimentation in the vision and language domains for multiple use cases of MTL. We find that our attacks can consistently achieve non-trivial success rates in terms of AUC and TPR, even when the adversary has no reference knowledge and chooses thresholds based on percentiles. Additionally, we attempt to understand the factors that lead to task-inference leakage by measuring how attack AUC varies for both the \textbf{strong} and \textbf{weak} adversaries as a function of the model's generalization gap and running ablation studies over different MTL hyperparameters, such as the embedding dimension and number of tasks.

\subsection{Potential Defenses}
% Group-level differential privacy
% Private MTL - Client Level Guarantee
% Clipping is not sufficient like in lira
One potential defense against our attack could be user-level (and group-level) differential privacy~\cite{dwork-2006-dmns}, where neighboring datasets are defined by the inclusion or omission of a user's entire contribution to the dataset, rather than an individual sample. While there are several works study user-level differential privacy when training machine learning models~\cite{levy-2021-learninguserdp, chua-2024-privacyunituserlevel, charles-2024-finetuningllmuserlevel} and estimating high dimensional means~\cite{cummings-2022-userlevelmean, agarwal-2025-personlevel}, only one work develops algorithms with \emph{client-level} privacy guarantees~\cite{hu2023privatemtl} when MTL is applied to collaborative learning. In contrast to observations for empirical defenses made in the membership-inference literature~\cite{carlini-2022-lira}, our attack succeeds even when the target MTL model is trained using gradient clipping and normalization. In fact, we use gradient clipping in our evaluation to ensure that the MTL models converge smoothly and do not overfit when learning over very few samples from each task.

% \subsection{Privacy Attacks on Groups}

